\chapter{Fetal Alcohol Spectrum Disorder}

\section*{Definition and Overview}
Fetal alcohol spectrum disorder (FASD) is the term for the range of lifelong physical, behavioural, and learning problems that can occur in a child whose mother drank alcohol during pregnancy. Alcohol passes freely through the placenta, and the developing brain is highly vulnerable to its effects at every stage of pregnancy; there is no known safe level of alcohol use while pregnant. FASD is a neurodevelopmental disability: it is not a behaviour problem or a sign of poor parenting, and it does not go away with age. However, with early diagnosis, understanding, and the right supports, children and adults with FASD can live well. The most severe form, fetal alcohol syndrome, includes characteristic facial features and growth failure, but most affected people have no distinctive facial appearance. This chapter explains the causes, features, diagnosis, and support for FASD.

\section*{Epidemiology}
FASD is the leading preventable cause of developmental disability in Australia. Conservative estimates suggest that at least 2\% of Australian children are born with FASD, meaning tens of thousands of people are affected, and rates are higher in communities with higher levels of drinking during pregnancy. Most affected people are never diagnosed, because there are no physical signs in the majority and diagnosis requires specialised multidisciplinary assessment. Internationally, prevalence estimates range from 1 to 5\% of the population. FASD is at least as common as autism spectrum disorder, yet it receives far less recognition, and the burden of unrecognised FASD falls heavily on individuals, families, schools, and the justice system.

\section*{Aetiology and Pathophysiology}
\begin{itemize}
    \item \textbf{Direct toxicity:} alcohol and its breakdown product acetaldehyde cross the placenta and damage developing brain cells, interfering with cell division, migration, and survival
    \item \textbf{Timing:} the first trimester carries the highest risk of structural malformations and facial changes, but brain damage can occur throughout pregnancy, including in the third trimester when the brain grows fastest
    \item \textbf{Dose and pattern:} binge drinking is particularly harmful, but regular drinking at low levels also carries risk; there is no threshold below which harm is excluded
    \item \textbf{Brain effects:} damage to the corpus callosum, hippocampus, and frontal lobes produces problems with learning, memory, attention, impulse control, and planning
    \item \textbf{Susceptibility:} maternal age, nutrition, genetics, and other substance use modify the effect, which is why outcomes differ between pregnancies with similar exposure
    \item \textbf{Physical effects:} disrupted development can cause growth restriction, facial dysmorphism, and malformations of the heart, kidneys, and bones
    \item \textbf{Lifelong condition:} the brain changes are permanent, although their functional impact can be greatly modified by environment and support
\end{itemize}

\section*{Clinical Features}
The features span brain-based difficulties and physical signs.
\begin{itemize}
    \item Cognitive impairments: difficulties with learning, memory, reasoning, and problem solving
    \item Attention and hyperactivity: distractibility, impulsivity, and difficulty regulating behaviour
    \item Executive function: trouble planning, organising, following multi-step instructions, and understanding consequences
    \item Language and social skills: delayed speech, difficulty reading social cues, and trouble making and keeping friends
    \item Daily living: difficulty managing money, time, and routines, and vulnerability to exploitation
    \item Physical signs in fetal alcohol syndrome only: a smooth philtrum, thin upper lip, and small eye openings, plus growth deficiency
    \item Associated problems: heart defects, hearing and vision problems, poor coordination, and sleep difficulties
    \item Mental health: anxiety, depression, and behavioural disorders are very common, especially without diagnosis and support
\end{itemize}

\section*{Differential Diagnosis}
The brain-based features overlap with many conditions, and careful assessment is needed.
\begin{itemize}
    \item Attention deficit hyperactivity disorder (ADHD)
    \item Autism spectrum disorder
    \item Intellectual disability and specific learning disorders
    \item The effects of trauma, neglect, and unstable caregiving
    \item Reactive attachment disorder
    \item Other genetic syndromes, such as Williams, Noonan, and fragile X syndromes
    \item Prematurity and other prenatal exposures, including drugs other than alcohol
\end{itemize}

\section*{When to Seek Medical Care}
Parents and carers should seek assessment if a child has known alcohol exposure in pregnancy or shows persistent difficulties with learning, behaviour, growth, or development. Early referral to a specialised FASD diagnostic service opens the door to funding, therapy, and school support.

\textbf{Call 000 (or local emergency services) for:}
\begin{itemize}
    \item Any medical emergency in a child
    \item Suicidal behaviour or statements, which are more common in older children and adults with FASD
    \item Any situation in which a child or young person is at risk of harm
\end{itemize}

\section*{Diagnosis and Investigations}
\begin{itemize}
    \item \textbf{Confirmed prenatal alcohol exposure:} a reliable history of drinking in pregnancy, from the mother or from records
    \item \textbf{Neurodevelopmental assessment:} comprehensive psychological testing of cognition, memory, attention, executive function, language, and motor skills
    \item \textbf{Physical examination:} growth measurement and examination for facial features and birth defects
    \item \textbf{Medical assessment:} hearing, vision, cardiac, and other reviews by a paediatrician
    \item \textbf{Diagnostic criteria:} Australian guidelines require severe impairment in at least three brain domains, with or without the facial features
    \item \textbf{Reports from home and school:} information about daily functioning from parents, teachers, and carers is essential
    \item \textbf{Genetic testing:} sometimes used to exclude other conditions
\end{itemize}

\section*{Management}
\begin{itemize}
    \item \textbf{Early intervention:} speech therapy, occupational therapy, physiotherapy, and preschool programs
    \item \textbf{Educational support:} individualised plans, routines, visual schedules, and teaching methods suited to the child's profile
    \item \textbf{Behaviour support:} consistent routines, simple instructions, positive reinforcement, and appropriate supervision, especially around money, safety, and social media
    \item \textbf{Medication:} treatment of co-occurring ADHD, anxiety, or depression where diagnosed
    \item \textbf{Family support and training:} helping caregivers understand that difficult behaviour reflects brain injury, not defiance
    \item \textbf{Transition support:} assistance moving from school into vocational training, work, and supported living
    \item \textbf{NDIS and advocacy:} funding for supports, and advocacy to protect vulnerable young people
    \item \textbf{Adult care:} mental health care, addiction services, and supported accommodation where needed
\end{itemize}

\section*{Complications}
\begin{itemize}
    \item Secondary disabilities when FASD is unrecognised: school failure, exclusion, contact with the justice system, substance misuse, and unsafe behaviour
    \item Mental health problems, including anxiety, depression, and elevated suicide risk
    \item Exploitation and victimisation from impaired judgement and suggestibility
    \item Difficulty with employment, relationships, and independent living
    \item Chronic health problems from birth defects
    \item High caregiver burden and family stress without adequate support
\end{itemize}

\section*{Prognosis and Outcomes}
FASD is lifelong, but outcomes are strongly shaped by early diagnosis and support. Children diagnosed young and raised in stable, understanding environments often finish school, find work, and live satisfying lives; protective factors include a consistent home, caregiver knowledge, and services that build on strengths. Without recognition, the same children are far more likely to experience school exclusion, mental illness, and justice involvement. There is no cure, and the most powerful intervention is prevention: not drinking alcohol in pregnancy. Research into supports that build executive function and self-regulation continues to improve prospects for those affected.

\section*{Prevention and Self-Care}
\begin{itemize}
    \item Do not drink alcohol if pregnant, trying to conceive, or not using effective contraception
    \item Stop drinking as soon as a pregnancy is discovered; stopping at any stage reduces harm
    \item Seek help from a doctor or alcohol support service if reducing drinking is difficult
    \item Support alcohol-free pregnancies through partners, families, and communities
    \item For parents of a child with FASD: build routines, simplify instructions, supervise high-risk situations, and connect with support groups and respite
    \item Advocate at school and with services, and seek NDIS assessment early
    \item Care for yourself as a caregiver, including counselling and peer networks
\end{itemize}

\section*{Sources and Further Reading}
The following reputable sources provide further information for healthcare students and patients seeking reliable, current advice about fetal alcohol spectrum disorder.
\begin{itemize}
    \item NOFASD Australia. \textit{FASD information and support.} https://www.nofasd.org.au/
    \item Healthdirect Australia. \textit{Alcohol and pregnancy.} https://www.healthdirect.gov.au/
    \item Telethon Kids Institute. \textit{FASD research and guidelines.} https://www.telethonkids.org.au/
    \item National Health and Medical Research Council. \textit{Alcohol and pregnancy guidelines.} https://www.nhmrc.gov.au/
    \item NHS. \textit{Fetal alcohol syndrome.} https://www.nhs.uk/conditions/fetal-alcohol-syndrome/
    \item World Health Organization. \textit{Alcohol and pregnancy.} https://www.who.int/
\end{itemize}
